
    %%%%%%%%%%%%%%%%%%%%%%%%%%%%%%%%%%%%%%%%%%%%%%%%%%%%%%%%%%%%%%%%%%%%%%%%%%%%%%%%
    %%% Classe do documento
    \documentclass[12pt, addpoints]{exam}
    \UseRawInputEncoding

    %%%%%%%%%%%%%%%%%%%%%%%%%%%%%%%%%%%%%%%%%%%%%%%%%%%%%%%%%%%%%%%%%%%%%%%%%%%%%%%%
    % preambulo.tex
    %
    % Arquivo de configuração de packages para uso o LaTeX, conforme minhas
    % preferências e modelos pessoais.
    %
    % Para maiores informações, visite:
    %    https://github.com/abrantesasf/latex
    %
    % NÃO ALTERE SE NÃO SOUBER O QUE ESTÁ FAZENDO!
    %%%%%%%%%%%%%%%%%%%%%%%%%%%%%%%%%%%%%%%%%%%%%%%%%%%%%%%%%%%%%%%%%%%%%%%%%%%%%%%%


    %%%%%%%%%%%%%%%%%%%%%%%%%%%%%%%%%%%%%%%%%%%%%%%%%%%%%%%%%%%%%%%%%%%%%%%%%%%%%%%%
    %%% Estruturas de controle:
    \input{static/utils/controles}

    %%%%%%%%%%%%%%%%%%%%%%%%%%%%%%%%%%%%%%%%%%%%%%%%%%%%%%%%%%%%%%%%%%%%%%%%%%%%%%%%
    %%% Compilação condicional em PDF:
    \input{static/utils/pdf}

    %%%%%%%%%%%%%%%%%%%%%%%%%%%%%%%%%%%%%%%%%%%%%%%%%%%%%%%%%%%%%%%%%%%%%%%%%%%%%%%%
    %%% Configurações de layout da página:
    \input{static/utils/layout}

    %%%%%%%%%%%%%%%%%%%%%%%%%%%%%%%%%%%%%%%%%%%%%%%%%%%%%%%%%%%%%%%%%%%%%%%%%%%%%%%%
    %%% Configurações para autores:
    \input{static/utils/autores}

    %%%%%%%%%%%%%%%%%%%%%%%%%%%%%%%%%%%%%%%%%%%%%%%%%%%%%%%%%%%%%%%%%%%%%%%%%%%%%%%%
    %%% Cabeçalho e rodapé:
    %%% Atenção! É MUITO DIFÍCIL ajustar apropriadamente os running headers
    %%% e running footers da primeira vez. Você pode confiar no LaTeX e deixar que
    %%% ele faça o ajuste automaticamente ou pode quebrar a cabeça e
    %%% tentar melhorar algo que o LaTeX já faz muito bem, mesmo que não fique
    %%% exatamente do jeito que você quer. O que você prefere? Se quiser ajustar
    %%% manualmente, descomente a linha abaixo e faça as configurações no arquivo
    %%% "utils/headings.tex".
    %\input{static/utils/headings}

    %%%%%%%%%%%%%%%%%%%%%%%%%%%%%%%%%%%%%%%%%%%%%%%%%%%%%%%%%%%%%%%%%%%%%%%%%%%%%%%%
    %%% Configuração para as fontes do documento:
    %%% Se você quiser usar fontes próprias ou do sistema, precisará ajustar as
    %%% configurações no arquivo "utils/fontes.tex".
    %%% ATENÇÃO: por padrão eu utilizo XeLaTeX com fontes proprietárias projetadas
    %%% por Matthew Butterick, adquiridas em https://mbtype.com, que não podem ser
    %%% distribuídas devido à licença de utilização. Se você usar XeLaTeX, você
    %%% DEVERÁ OBRIGATORIAMENTE ajustar as fontes no arquivo "utils/fontes.tex".
    \input{static/utils/fontes}

    %%%%%%%%%%%%%%%%%%%%%%%%%%%%%%%%%%%%%%%%%%%%%%%%%%%%%%%%%%%%%%%%%%%%%%%%%%%%%%%%
    %%% Sumário:
    \input{static/utils/toc}

    %%%%%%%%%%%%%%%%%%%%%%%%%%%%%%%%%%%%%%%%%%%%%%%%%%%%%%%%%%%%%%%%%%%%%%%%%%%%%%%%
    %%% Matemática:
    \input{static/utils/matematica}

    %%%%%%%%%%%%%%%%%%%%%%%%%%%%%%%%%%%%%%%%%%%%%%%%%%%%%%%%%%%%%%%%%%%%%%%%%%%%%%%%
    %%% Notas de rodapé:
    \input{static/utils/footnotes}

    %%%%%%%%%%%%%%%%%%%%%%%%%%%%%%%%%%%%%%%%%%%%%%%%%%%%%%%%%%%%%%%%%%%%%%%%%%%%%%%%
    %%% Referências cruzadas, links, citações:
    \input{static/utils/crossrefs}

    %%%%%%%%%%%%%%%%%%%%%%%%%%%%%%%%%%%%%%%%%%%%%%%%%%%%%%%%%%%%%%%%%%%%%%%%%%%%%%%%
    %%% Configurações para os hiperlinks em PDF:
    \input{static/utils/pdfconfig}

    %%%%%%%%%%%%%%%%%%%%%%%%%%%%%%%%%%%%%%%%%%%%%%%%%%%%%%%%%%%%%%%%%%%%%%%%%%%%%%%%
    %%% Glossário e índice remissivo:
    \input{static/utils/glossindex}

    %%%%%%%%%%%%%%%%%%%%%%%%%%%%%%%%%%%%%%%%%%%%%%%%%%%%%%%%%%%%%%%%%%%%%%%%%%%%%%%%
    %%% Referências bibliográficas:
    \input{static/utils/bibliografia}

    %%%%%%%%%%%%%%%%%%%%%%%%%%%%%%%%%%%%%%%%%%%%%%%%%%%%%%%%%%%%%%%%%%%%%%%%%%%%%%%%
    %%% Cores:
    \input{static/utils/cores}

    %%%%%%%%%%%%%%%%%%%%%%%%%%%%%%%%%%%%%%%%%%%%%%%%%%%%%%%%%%%%%%%%%%%%%%%%%%%%%%%%
    %%% Computação:
    \input{static/utils/computacao}

    %%%%%%%%%%%%%%%%%%%%%%%%%%%%%%%%%%%%%%%%%%%%%%%%%%%%%%%%%%%%%%%%%%%%%%%%%%%%%%%%
    %%% Gráficos:
    \input{static/utils/graficos}

    %%%%%%%%%%%%%%%%%%%%%%%%%%%%%%%%%%%%%%%%%%%%%%%%%%%%%%%%%%%%%%%%%%%%%%%%%%%%%%%%
    %%% Ícones:
    \input{static/utils/icones}

    %%%%%%%%%%%%%%%%%%%%%%%%%%%%%%%%%%%%%%%%%%%%%%%%%%%%%%%%%%%%%%%%%%%%%%%%%%%%%%%%
    %%% Blocos:
    \input{static/utils/blocos}

    %%%%%%%%%%%%%%%%%%%%%%%%%%%%%%%%%%%%%%%%%%%%%%%%%%%%%%%%%%%%%%%%%%%%%%%%%%%%%%%%
    %%% Boxes coloridos:
    \input{static/utils/colorbox}

    %%%%%%%%%%%%%%%%%%%%%%%%%%%%%%%%%%%%%%%%%%%%%%%%%%%%%%%%%%%%%%%%%%%%%%%%%%%%%%%%
    %%% Tabelas:
    \input{static/utils/tabelas}

    %%%%%%%%%%%%%%%%%%%%%%%%%%%%%%%%%%%%%%%%%%%%%%%%%%%%%%%%%%%%%%%%%%%%%%%%%%%%%%%%
    %%% Ambientes de listas:
    \input{static/utils/listas}

    %%%%%%%%%%%%%%%%%%%%%%%%%%%%%%%%%%%%%%%%%%%%%%%%%%%%%%%%%%%%%%%%%%%%%%%%%%%%%%%%
    %%% Ambientes floats e similares:
    \input{static/utils/floats}

    %%%%%%%%%%%%%%%%%%%%%%%%%%%%%%%%%%%%%%%%%%%%%%%%%%%%%%%%%%%%%%%%%%%%%%%%%%%%%%%%
    %%% Ajustes para a classe EXAM:
    \input{static/utils/exam}

    %%%%%%%%%%%%%%%%%%%%%%%%%%%%%%%%%%%%%%%%%%%%%%%%%%%%%%%%%%%%%%%%%%%%%%%%%%%%%%%%
    %%% Ajustes para textos:
    \input{static/utils/texto}

    %%%%%%%%%%%%%%%%%%%%%%%%%%%%%%%%%%%%%%%%%%%%%%%%%%%%%%%%%%%%%%%%%%%%%%%%%%%%%%%%
    %%% Ambientes, aliases, comandos e customizações em geral para serem
    %%% utilizadas nos documentos. Defina aqui qualquer customização específica
    %%% para seus documentos!
    \input{static/utils/customizacoes}

    %%%%%%%%%%%%%%%%%%%%%%%%%%%%%%%%%%%%%%%%%%%%%%%%%%%%%%%%%%%%%%%%%%%%%%%%%%%%%%%%
    %%% Ajuste do layout, espaçamento de linhas e etc.:
    \geometry{a4paper, portrait, left=2.0cm, right=2.0cm, top=2.0cm, bottom=2.0cm}
    \singlespacing

    %%%%%%%%%%%%%%%%%%%%%%%%%%%%%%%%%%%%%%%%%%%%%%%%%%%%%%%%%%%%%%%%%%%%%%%%%%%%%%%%
    %%% Configurações de cabeçalho e rodapé (não use fancyhdr pois ocorre conflito):
    \pagestyle{headandfoot}
    \runningheadrule
    \firstpageheader{}{}{}
    \runningheader{Sem disciplina}{}{Avaliação}
    \firstpagefooter{}{}{Página \thepage\ de \numpages}
    \runningfooter{}{}{Página \thepage\ de \numpages}
    \runningfootrule

    %%%%%%%%%%%%%%%%%%%%%%%%%%%%%%%%%%%%%%%%%%%%%%%%%%%%%%%%%%%%%%%%%%%%%%%%%%%%%%%%
    %%% Configurações para as propriedades do PDF:
    \hypersetup{
    hidelinks,           % Comente para web, descomente para impressão
    colorlinks=true,      % True para web, False para impressão
    pdftitle=AAAAA: Avaliação,
    pdfauthor=Matheus Endlich Silveira,
    pdfsubject=Sem disciplina,
    pdfkeywords=['Sem tag'],
    pdfinfo={
        CreationDate={}, % Ex.: D:AAAAMMDDHH24MISS
        ModDate={}       % Ex.: D:AAAAMMDDHH24MISS
    }
    }
    \input{static/utils/pdfconfig}

    %%%%%%%%%%%%%%%%%%%%%%%%%%%%%%%%%%%%%%%%%%%%%%%%%%%%%%%%%%%%%%%%%%%%%%%%%%%%%%%%
    %%% Começa o documento
    \begin{document}

    %%%%%%%%%%%%%%%%%%%%%%%%%%%%%%%%%%%%%%%%%%%%%%%%%%%%%%%%%%%%%%%%%%%%%%%%%%%%%%%%
    %%% Folha de instruções padronizadas da UVV para a prova
    \pdfbookmark[1]{Instruções}{instrucoes}
    \begin{figure}[H]
        \begin{center}
            \includegraphics[scale=0.3]{{static/imgs/logo-uvv.png}}
        \end{center}	
    \end{figure}

    \vspace{-1cm}

    \begin{table}[H]
        \centering
        \begin{tabular}{|llll|}
            \hline
            \multicolumn{2}{|l|}{\textbf{Disciplina:} Sem disciplina} & 
            \multicolumn{1}{l|}{\multirow{3}{*}{\textbf{\makecell{Nota:\\ \,\\ \,}}}} & 
            \multirow{3}{*}{\textbf{\makecell{Coordenador:\\ \,\\ \,}}} \\ 
            \cline{1-2}
            \multicolumn{2}{|l|}{\textbf{Professor:} Matheus Endlich Silveira \hspace{4.72cm} } & 
            \multicolumn{1}{l|}{} & \\ 
            \cline{1-2}
            \multicolumn{2}{|l|}{\textbf{Aluno:}} & 
            \multicolumn{1}{l|}{} & \\ 
            \hline
            \multicolumn{1}{|l|}{\textbf{Turma:} \hspace{6.3cm} } & 
            \multicolumn{1}{l|}{\textbf{Semestre:}} & 
            \multicolumn{2}{l|}{\textbf{Valor:} 10 pontos} \\ 
            \hline
            \multicolumn{1}{|l|}{\textbf{Data:}} & 
            \multicolumn{3}{l|}{\textbf{Avaliação:}} \\ 
            \hline
        \end{tabular}
    \end{table}

    \begin{center}
        \fbox{\setlength\fboxsep{0.3cm} \fbox{\parbox{15.4cm}{
            \begin{center}
                \textbf{INSTRUÇÕES PARA A REALIZAÇÃO DESTA AVALIAÇÃO:}
            \end{center}
            \begin{itemize}[leftmargin=*]
                \item Esta é a primeira avaliação da disciplina Sem disciplina, 
                    e engloba o conteúdo referente Sem tag.
                \item Esta avaliação contém 14 questões objetivas, \textbf{todas obrigatórias}.
                \item \textbf{Leia atentamente cada questão!} As questões objetivas terão uma,
                    e apenas uma única, resposta a ser marcada.
                \item \textbf{Desligue o celular} antes de começar. A avaliação é
                    \textbf{individual} e \textbf{sem consulta}.
                \item As questões podem ser respondidas, na prova, com lápis ou caneta. Ao final
                    da prova \textbf{{VOCÊ DEVERÁ PREENCHER O GABARITO COM CANETA
                    DE COR PRETA}} \textbf{\underline{OU AZUL}} para que seja feita a correção
                    automatizada através da leitura do padrão de respostas marcado no
                    gabarito.
                \item \textbf{CUIDADO ao preencher o gabarito} pois eles são \textbf{INDIVIDUAIS
                    e NOMEADOS} (cada estudante tem um gabarito próprio específico, com um
                    código de identificação único). \textbf{Questões rasuradas são anuladas}
                    pelo software de leitura do gabarito.
                \item Ao final da prova, \textbf{DEVOLVA A PROVA E O GABARITO} para o professor.
                \item Siga todas as normas de \textbf{Integridade Acadêmica} da disciplina.
                    Alunos que forem flagrados com qualquer espécie de ``cola'' ou trocando
                    informações com outros alunos terão suas avaliações recolhidas, as notas
                    zeradas, e a situação será encaminhada para a coordenação para a aplicação
                    das penalidades previstas pela universidade.
                \item Com 90 minutos de avaliação você terá tempo de até 6,min por questão,
                    em média.
                \item Boa avaliação!
            \end{itemize}
            \vspace{2cm}
        }}}
    \end{center}
    \newpage

    %%%%%%%%%%%%%%%%%%%%%%%%%%%%%%%%%%%%%%%%%%%%%%%%%%%%%%%%%%%%%%%%%%%%%%%%%%%%%%%%
    %%% Inicia o bloco de questões:
    \begin{questions}

    \newpage
    \question

O usuário que não é da área de tecnologia, costuma ser de alta gerência, que é

especialista no negócio e tem a responsabilidade central pelas políticas de

dados da empresa é o:
\begin{choices}
\choice Gerente de Tecnologia da Informação (ITM: IT Manager)
\choice Coordenador de Tecnologia da Informação (ITS: IT Supervisor)
\choice Administrador do Banco de Dados (DBA: Database Administrator)
\choice Administrador de Dados (DA: Data Administrator)
\choice Encarregado da Proteção de Dados (DPO: Data Protection Officer)
\end{choices}

\question

Você estava ouvindo a discussão de dois colegas, o João e a Maria. Enquanto o

João afirmava que o conceito de banco de dados (BD) é extremamente diferente do

conceito de sistema de gerenciamento de bancos de dados (SGBD), a Maria

discordava e afirmava que o banco de dados nada mais é do que um componente do

SGBD. Nesse momento o professor entrou na sala, percebeu a discussão e falou

que:
\begin{choices}
\choice Os dois estavam errados pois não são conceitos separados nem componentes, são conceitos sinônimos.
\choice Os dois estavam certos pois o conceito, na prática, depende da situação específica que se está considerando.
\choice O João estava certo já que o banco de dados é, realmente, um conceito diferente do conceito do sistema de gerenciamento de bancos de dados, e existe de forma independente.
\choice Não é possível saber quem estava certo ou errado pois não definiram qual o modelo de dados que estava sendo considerado na discusão.
\choice A Maria estava certa já que o banco de dados é, realmente, um componente interno do sistema de gerenciamento de bancos de dados e não existe de forma independente.
\end{choices}

\question

Pode-se afirmar que o gráfico da função \(y = 2 + \frac{1}{x - 1}\) é o gráfico da função \(y = \frac{1}{x}\):
\begin{choices}
\choice transladado uma unidade para a esquerda e duas unidades para cima
\choice transladado uma unidade para a esquerda e duas unidades para baixo
\choice transladado uma unidade para a direita e duas unidades para baixo
\choice nenhuma das anteriores.
\choice transladado uma unidade para a direita e duas unidades para cima
\end{choices}

\question

Em um banco de dados relacional, o que é uma chave primária?


\begin{choices}
\choice Nenhuma das respostas acima
\choice É um atributo (ou um conjunto de atributos) que permite identificar única e exclusivamente cada registro da tabela
\choice É uma validação que evita o cadastro de dados errados
\choice É qualquer coluna ou conjunto de colunas que podem servir como índice
\choice É denominada por FK
\end{choices}

\question

(Adaptado do ENADE 2005) Analise o relacionamento existente entre as tabelas

``clubes' e ``jogadores', nas alternativas abaixo, e marque a alternativa que

indica que todo jogador deve pertencer a um único clube, e todo clube poder ter

ou não jogadores.
\begin{choices}
\choice \includegraphics[width=0.5\linewidth]{static/imgs/defaul.png} \vspace{0.5cm}
\choice \includegraphics[width=0.5\linewidth]{static/imgs/defaul.png} \vspace{0.5cm}
\choice \includegraphics[width=0.5\linewidth]{static/imgs/defaul.png} \vspace{0.5cm}
\choice \includegraphics[width=0.5\linewidth]{static/imgs/defaul.png} \vspace{0.5cm}
\choice \includegraphics[width=0.5\linewidth]{static/imgs/defaul.png} \vspace{0.5cm}
\end{choices}

\question

Quando nos referimos à propriedade auto-descritiva dos bancos de dados estamos

trabalhando, fundamentalmente, com o conceito de:
\begin{choices}
\choice Armazenamento
\choice Usuários
\choice Relacionamentos
\choice Metadados
\choice SQL
\end{choices}

\question

Qual o valor do atributo \( E.val \) após a análise da expressão "4 / 2 / 2" para o esquema de tradução a seguir?



\begin{itemize}

    \item \( E \rightarrow T / E1 \ \{ E.val = T.val / E1.val \} \)

    \item \( E \rightarrow T \ \{ E.val = T.val \} \)

    \item \( T \rightarrow \text{digito} \ \{ T.val = \text{val(digito)} \} \)

\end{itemize}


\begin{choices}
\choice 2
\choice 3
\choice 1
\choice 8
\choice 4
\end{choices}

\question

O contador da figura abaixo é:

\begin{figure}

    \centering

    \includegraphics[width=0.5\linewidth]{static/imgs/defaul.png}

    \label{fig:enter-label}

\end{figure}
\begin{choices}
\choice síncrono
\choice anisócrono
\choice auto-sincronizado
\choice isócrono
\choice assíncrono
\end{choices}

\question

\textbf Supondo a Relação \texttt{PROJ (PNO, Nome, Orçam)}, com chave primária \texttt{PNO} e a Relação \texttt{DSG (ENO, PNO, Dur, Resp)}, com chave primária \{\texttt{ENO, PNO}\} e chave estrangeira \texttt{PNO} em relação a \texttt{PROJ}, a asserção abaixo \textbf{NÃO} expressa:



\begin{equation}\forall g \in \texttt{DSG}, \exists j \in \texttt{PROJ} : g.\texttt{PNO} = j.\texttt{PNO}\end{equation}
\begin{choices}
\choice Uma restrição a ser verificada na atualização de tuplas em $\texttt{DSG}$.
\choice Uma restrição que define um estado consistente do banco de dados.
\choice Uma restrição a ser verificada na inserção de tuplas em $\texttt{DSG}$.
\choice Uma restrição de integridade de chave primária em $\texttt{PROJ}$.
\choice Uma restrição de integridade de chave estrangeira em $\texttt{DSG}$.
\end{choices}

\question

Em SQL, qual comando é utilizado para combinar linhas de duas ou mais tabelas?


\begin{choices}
\choice JOIN
\choice COMBINE
\choice MERGE
\choice CONNECT
\choice LINK
\end{choices}

\question

O Modelo Espiral de desenvolvimento de software, proposto por Barry Boehm na década de 1970, apresenta as seguintes afirmações:

\begin{enumerate}[label=\Roman*.]

    \item Suas atividades de desenvolvimento são conduzidas por riscos;

    \item Cada ciclo da espiral inclui 4 passos: passo 1 - identificação dos objetivos; passo 2 - avaliação das alternativas tendo em vista os objetivos e os riscos (incertezas, restrições) do desenvolvimento; passo 3 - desenvolvimento de estratégias (simulação, prototipagem) para resolver riscos; e passo 4 - planejamento do próximo passo e continuidade do processo determinada pelos riscos restantes;

    \item É um modelo evolutivo em que cada passo pode ser representado por um quadrante num diagrama cartesiano: assim, na dimensão radial da espiral tem-se o custo acumulado dos vários passos do desenvolvimento, enquanto na dimensão angular tem-se o progresso do projeto.

\end{enumerate}

Levando-se em conta as três afirmações I, II e III acima, identifique a única alternativa válida:


\begin{choices}
\choice Apenas a I e a III estão corretas;
\choice Apenas a II e a III estão corretas;
\choice As afirmações I, II e III estão corretas;
\choice Apenas a III está correta.
\choice Apenas a I e a II estão corretas;
\end{choices}

\question

No que diz respeito as vantagens da arquitetura de micro-núcleo para sistemas operacionais em relação a arquiteturas de núcleo monolítico, quais das seguintes afirmações são verdadeiras?



\begin{itemize}

    \item[I] A arquitetura de micro-núcleo facilita a depuração do SO.

    \item[II] A arquitetura de micro-núcleo permite um número menor de mudanças de contexto.

    \item[III] A arquitetura de micro-núcleo facilita a reconfiguração de serviços do SO pois a maioria deles reside em espaço de usuário.

\end{itemize}
\begin{choices}
\choice II e III
\choice I e III
\choice I e II
\choice Todas são verdadeiras
\choice Apenas I
\end{choices}

\question

Considere um sistema distribuído onde cada nó precisa obter um bloqueio (\textit{lock}) antes de acessar qualquer serviço no sistema. Qual das estratégias a seguir \textbf{não} seria eficaz para evitar impasses (\textit{deadlocks})?
\begin{choices}
\choice Fazer com que cada nó reinicie sua execução se um pedido de bloqueio não é concedido após um longo tempo de espera. O pedido de bloqueio é re-enviado após um tempo aleatório
\choice Forçar cada nó a obter todos os bloqueios de que necessita no início de sua execução e reiniciar a execução se algum bloqueio não é concedido
\choice Instalar um serviço de detecção de impasses no sistema distribuído e reiniciar os nós que atinjam um impasse
\choice Associar prioridades aos nós e criar filas de prioridades para cada serviço
\choice Numerar os serviços e exigir que cada nó solicite os bloqueios dos serviços em ordem crescente
\end{choices}

\question

Para a gramática a seguir, qual o conjunto de terminais que pode aparecer como primeiro terminal após o não-terminal \( A \), em qualquer forma sentencial gerada pela gramática abaixo (isto é, não necessariamente imediatamente após \( A \)), onde \( \epsilon \) representa a sentença vazia?



$S \rightarrow ABCDd$\\

$A \rightarrow aA \mid \epsilon$\\

$B \rightarrow bC \mid \epsilon$\\

$C \rightarrow cD \mid \epsilon$\\

$D \rightarrow e$\\
\begin{choices}
\choice \{d\}
\choice \{b, c, e\}
\choice \{b, c, d, e\}
\choice \{e\}
\choice \{b\}
\end{choices}



    %%%%%%%%%%%%%%%%%%%%%%%%%%%%%%%%%%%%%%%%%%%%%%%%%%%%%%%%%%%%%%%%%%%%%%%%%%%%%%%%
    %%% Finaliza o bloco de questões:
    \end{questions}
    \end{document}
    